\section{Related work}\label{sec:related}

\paragraph{Content-addressed code.}
Naming data by the hash of its content goes back to Merkle's hash trees~\cite{merkle1987}, and
underlies version control, the Nix package manager~\cite{dolstra2006thesis} and content-addressed
storage networks~\cite{benet2014ipfs}. Unison~\cite{unison} names each definition by a hash of
its syntax tree with names abstracted, and hashes mutually recursive definitions as a cycle;
\VG{} borrows this treatment of cycles for its future recursive nodes. Two differences matter.
\VG{} hashes the contract together with the body, so that a hash names a function \emph{and}
what is proven about it, and every proof is cached by the hash of its query; and it does not
expect different implementations to meet at one hash, using spec classes instead
(\cref{thm:conp}). Hashing modulo $\alpha$-equivalence~\cite{maziarz2021hashing} is how its
contract hash ignores bound names.

\paragraph{Term graphs, hash-consing and canonical forms.}
\VG{} bodies are term graphs~\cite{ariola1996equational}, and their canonical form is the fully
shared graph of their unraveling, computed by hash-consing~\cite{filliatre2006hashcons}.
Single-assignment form~\cite{cytron1991ssa} and the sea-of-nodes
representation~\cite{click1995simple} also represent programs as dataflow graphs, as do the
dialects of MLIR~\cite{lattner2021mlir}; they serve optimization, not identity, and do not carry
contracts. Equality saturation~\cite{tate2009equality,willsey2021egg,zhang2023egglog} computes
larger equivalences than commutativity and could make more programs share a canonical form;
\cref{thm:conp} bounds what any such efficient procedure can achieve. The soundness of merging
pure ops rests on the cartesian structure of pure computation~\cite{fox1976coalgebras}.

\paragraph{Verification-aware languages.}
\VG's contracts follow Hoare logic~\cite{hoare1969,floyd1967} and design by
contract~\cite{meyer1992}; modular verification against callee contracts is the discipline of
Boogie~\cite{barnett2005boogie}, Dafny~\cite{leino2010dafny}, Why3~\cite{filliatre2013why3},
F*~\cite{swamy2016fstar}, Liquid Haskell~\cite{vazou2014refinement} and
Verus~\cite{lattuada2023verus}. Those systems verify text written for people and accept
quantified specifications that need triggers or user lemmas. \VG{} restricts loops to bounded
combinators with mandatory invariants, so that every obligation is one quantifier-free query
accepted by Z3~\cite{demoura2008z3}, cvc5~\cite{barbosa2022cvc5} and
Bitwuzla~\cite{niemetz2023bitwuzla}, and returns every judgment as data. Its rule for matching
registry candidates is behavioral subtyping~\cite{liskov1994behavioral}. The SMT-LIB theory of
floating point~\cite{brain2015fp,smtlib} fixes the one-NaN domain that \cref{sec:float} adopts
for the language itself.

\paragraph{Effects and capabilities.}
Effect systems date back to Lucassen and Gifford~\cite{lucassen1988polymorphic}; monads and
premonoidal categories give effectful computation its semantics~\cite{moggi1991notions,
power1997premonoidal}, and string diagrams for effectful categories draw effects on a threaded
wire~\cite{roman2022promonads}, which is the role of \VG's effect token. Linear
types~\cite{wadler1990linear} explain why the token must be used exactly once
(\cref{lem:chain,lem:eff-never-merged}). Capabilities as unforgeable
references~\cite{dennis1966programming} and object-capability
security~\cite{miller2006robust} are the model of \VG's handles (\cref{prop:caps}).

\paragraph{Floating-point reasoning.}
The pitfalls of verifying floating-point programs are well documented~\cite{goldberg1991,
monniaux2008pitfalls}; Flocq~\cite{boldo2011flocq} formalizes IEEE arithmetic in Coq, and the
handbook of Muller et al.~\cite{muller2018handbook} collects the lemmas we use, among them
Sterbenz's~\cite{sterbenz1974}. Our case-study proof is a pen-and-paper instance of such
reasoning; mechanizing it, or obtaining it from the solver's floating-point theory, is what the
verifier will do.

\paragraph{Proof-carrying code and certified compilation.}
Proof-carrying code~\cite{necula1997pcc} ships proofs with code so that a consumer can check
them; CompCert~\cite{leroy2009formal} and CakeML~\cite{kumar2014cakeml} verify the compiler
itself. \VG{} neither ships proofs inside nodes nor verifies its compiler: proofs are stored by
query hash and re-checked, the compiler is in the trusted base, and miscompilation is attacked by
differential testing and translation validation~\cite{pnueli1998translation}. Alive2's bounded
translation validation~\cite{lopes2021alive2} is the planned check for lifted legacy code.

\paragraph{Language models, verifiers and constrained decoding.}
Language models increasingly write verified code: on a benchmark of 12,504 formal
specifications, off-the-shelf models succeeded on 82\% of the Dafny tasks, 44\% of the Verus
tasks and 27\% of the Lean tasks~\cite{bursuc2025vericoding},
and Dafny has been proposed as a verified intermediate language between a request and the
shipped code~\cite{li2025dafny}, the architecture \VG{} adopts with a graph instead of text.
Constrained decoding enforces syntax and types during generation~\cite{poesia2022synchromesh,
mundler2025type}; it distorts the model's distribution~\cite{park2024grammar}, which sequential
Monte Carlo can correct~\cite{loula2025syntactic}. \VG's define-before-use rule makes every
prefix checkable, so type-level masks can be exact; its hosted mode uses tools instead, because
hosted APIs expose no per-token masking. Agent performance depends on the interface the agent
is given~\cite{yang2024sweagent}; in a study of regular-expression synthesis, counterexample
feedback let an agent solve 90\% of 30 tasks within four turns, against 23\% with error-only
feedback~\cite{badrinarayan2026counterexamples}. This is why every \VG{} failure leads with a
replayable counterexample. Whether a formal specification
states what a user meant cannot be decided by any algorithm~\cite{lahiri2024evaluating}; \VG{}
renders every contract in plain words for confirmation and leaves intent validation to research.

\paragraph{Transparency and supply chain.}
The registry's ingestion log follows Certificate Transparency~\cite{rfc6962} and tamper-evident
logging~\cite{crosby2009efficient}; software signing and provenance systems such as
Sigstore~\cite{newman2022sigstore} and in-toto~\cite{torresarias2019intoto} address the same
attestations. Self-hosting is planned with the stage-2/stage-3 fixpoint and diverse
double-compiling~\cite{wheeler2005countering}, the standard answer to Thompson's
attack~\cite{thompson1984reflections}. Checkable SMT certificates in the Alethe
format~\cite{schurr2021alethe}, checked by Carcara~\cite{andreotti2023carcara}, are the planned
successor of today's re-solving with pinned solvers.
